\documentclass[11pt]{article}
\usepackage[margin=1in]{geometry}
\usepackage{amsmath,amssymb,amsthm}
\usepackage{enumitem}
\usepackage[T1]{fontenc}
\usepackage{lmodern}
\usepackage{hyperref}

\newtheorem{theorem}{Theorem}
\newtheorem{lemma}{Lemma}
\theoremstyle{definition}
\newtheorem{definition}{Definition}
\newtheorem{remark}{Remark}

\newcommand{\N}{\mathbb N}
\newcommand{\calD}{\mathcal D}
\newcommand{\calG}{\mathcal G}
\newcommand{\supp}{\operatorname{supp}}

\title{A harmonic counterexample for the power-profile \(f\)-greedy question}
\author{Run \texttt{fa\_banach\_001}, lane 7}
\date{2026-06-21}

\begin{document}
\maketitle

\section*{Source question}

The source paper is Pablo M. Berna and Hung Viet Chu,
\emph{On Some Characterizations of Greedy-type Bases}, arXiv:2201.12029.
For \(f:\N\to\mathbb R\) and a finite set
\(A=\{a_1<\cdots<a_m\}\), define
\[
        1_{f,A}=\sum_{r=1}^m f(r)e_{a_r},
        \qquad
        \calD_m^f(x)=
        \inf_{\substack{|A|=m\\ \alpha\in\mathbb R}}
        \|x-\alpha 1_{f,A}\|.
\]
A basis is called \(f\)-greedy if there is \(C_f\) such that
\[
        \|x-\calG_m(x)\|\le C_f \calD_m^f(x)
        \qquad (x\in X,\ m\in\N).
        \tag{1}
\]
The first future-research question in the source asks for a characterization
of the functions \(f\) for which (1) is equivalent to ordinary greediness.
The source proves the equivalence for regular \(f\), constructs non-greedy
examples for sparse \(f\), and explicitly points to
\[
        f(n)=n^{-p},\qquad p>0,
\]
as a class outside both results.

\section*{Counterexample theorem}

\begin{theorem}\label{thm:counterexample}
For every \(p>0\), there is a normalized Schauder basis which is
\(f_p\)-greedy for \(f_p(n)=n^{-p}\) but is not greedy.  In fact the same
basis works for every \(p>0\).
\end{theorem}

\section*{Idea of the proof}

Use the harmonic norm
\[
        \|x\|_h=\max\left\{\|x\|_\infty,\sum_{n=1}^\infty \frac{|x_n|}{n}\right\}.
\]
Initial intervals have norm \(1+1/2+\cdots+1/m\), while \(m\) coordinates far
out have norm \(1\), so the unit vector basis is not democratic.  The reason
it is still \(n^{-p}\)-greedy is that the weighted tail
\[
        \sum_{r\ge q}\frac{r^{-p}}{r}
\]
is comparable to \(q^{-p}\).  Thus the first missed \(f_p\)-coefficient pays
for the whole remaining harmonic tail.

\section*{Proof}

Let
\[
        X=\left\{x=(x_n)\in c_0:
        \sum_{n=1}^\infty \frac{|x_n|}{n}<\infty\right\}
\]
with norm
\[
        \|x\|=\max\left\{\sup_n |x_n|,
        \sum_{n=1}^\infty \frac{|x_n|}{n}\right\}.
        \tag{2}
\]
The canonical unit vectors \((e_n)\) form a normalized, \(1\)-unconditional
Schauder basis.  Indeed, the coordinate projections are contractive, and
partial sums converge in both the \(c_0\) and weighted \(\ell_1\) parts of
(2).

The basis is not democratic.  If \(A_m=\{1,\ldots,m\}\), then
\[
        \|1_{A_m}\|=\sum_{n=1}^m \frac{1}{n}.
\]
On the other hand, for each \(m\) one may choose \(N\) so large that
\(\sum_{n=N+1}^{N+m}1/n\le 1\).  Then
\[
        \left\|\sum_{n=N+1}^{N+m}e_n\right\|=1.
\]
The democracy constants therefore tend to infinity.  Since the basis is
unconditional, the Konyagin--Temlyakov characterization implies that it is not
greedy.

It remains to prove \(f_p\)-greediness.  Fix \(p>0\) and write
\[
        f(r)=r^{-p}.
\]
Let \(x\in X\), let \(\Lambda\) be the natural order-\(m\) greedy set of \(x\),
and put \(y=x-P_\Lambda x\).  We shall show that for every \(m\)-point set
\(A=\{a_1<\cdots<a_m\}\) and every scalar \(\alpha\),
\[
        \|y\|\le \left(3+\frac{2}{p}\right)
        \|x-\alpha 1_{f,A}\|.
        \tag{3}
\]
Taking the infimum over \(A\) and \(\alpha\) proves (1).

Set
\[
        z=x-\alpha 1_{f,A},\qquad \delta=\|z\|.
\]
First, \(\|y\|_\infty\le \delta\).  Indeed, let \(i\notin\Lambda\).  If
\(i\notin A\), then \(y_i=x_i=z_i\).  If \(i\in A\setminus\Lambda\), then
because \(|A|=|\Lambda|\), there is \(j\in\Lambda\setminus A\).  Since
\(\Lambda\) is greedy, \(|x_j|\ge |x_i|\), and \(z_j=x_j\); hence
\(|y_i|=|x_i|\le\delta\).

For the harmonic part, write
\[
        R=\{r\in\{1,\ldots,m\}: a_r\notin\Lambda\}.
\]
Then
\begin{equation}
\begin{aligned}
        \sum_{i\notin\Lambda}\frac{|x_i|}{i}
        &\le
        \sum_{i\notin\Lambda}\frac{|z_i|}{i}
        +|\alpha|\sum_{r\in R}\frac{r^{-p}}{a_r}  \\
        &\le \delta+
        |\alpha|\sum_{r\in R}\frac{r^{-p}}{a_r}.
\end{aligned}
\tag{4}
\end{equation}
If \(R=\varnothing\), (4) already gives the desired estimate.  Otherwise let
\(q=\min R\).  Since \(a_q\in A\setminus\Lambda\), choose
\(j\in\Lambda\setminus A\).  Then \(|x_j|\ge |x_{a_q}|\), while
\[
        z_{a_q}=x_{a_q}-\alpha q^{-p},
        \qquad z_j=x_j.
\]
Consequently
\[
        \delta\ge
        \max\{|x_{a_q}|,\ |x_{a_q}-\alpha q^{-p}|\}
        \ge \frac{|\alpha|q^{-p}}{2}.
        \tag{5}
\]
Also \(a_r\ge r\) for every \(r\), so
\[
        \sum_{r\in R}\frac{r^{-p}}{a_r}
        \le
        \sum_{r=q}^\infty \frac{1}{r^{p+1}}
        \le \left(1+\frac{1}{p}\right)q^{-p}.
        \tag{6}
\]
Combining (5) and (6) gives
\[
        |\alpha|\sum_{r\in R}\frac{r^{-p}}{a_r}
        \le 2\left(1+\frac{1}{p}\right)\delta.
\]
Together with (4),
\[
        \sum_{i\notin\Lambda}\frac{|x_i|}{i}
        \le \left(3+\frac{2}{p}\right)\delta.
\]
The supremum estimate and the weighted estimate prove (3), so the basis is
\(f_p\)-greedy.  This completes the counterexample.

\section*{Verification notes}

The proof is direct and does not use computation.  The only source-paper input
needed for context is the definition of \(f\)-greediness and the
Konyagin--Temlyakov characterization of greedy bases as unconditional plus
democratic.

The earlier partial packet
\[
\texttt{solutions/partial/2201.12029\_power\_profile\_log\_democracy/}
\]
is consistent with this counterexample: the present basis has precisely the
harmonic logarithmic democracy growth that the partial theorem allowed.

\section*{Novelty check}

The run indexes were searched for arXiv:2201.12029, \(f\)-greedy,
\(\calD_m^f\), power-profile terminology, \(1/n^p\), and the earlier partial
packet.  Web searches on 2026-06-21 for exact phrases around
``f-greedy \(n^{-p}\) counterexample'', ``\(D_m^f\) harmonic greedy basis'',
and the source title with ``\(1/n^p\)'' found no later matching answer.

\section*{References}

\begin{enumerate}[label={[\arabic*]}]
\item P. M. Berna and H. V. Chu,
      \emph{On Some Characterizations of Greedy-type Bases},
      arXiv:2201.12029.
\item S. V. Konyagin and V. N. Temlyakov,
      A remark on greedy approximation in Banach spaces,
      East J. Approx. 5 (1999), 365--379.
\end{enumerate}

\end{document}
